\setcounter{chapter}{4}
\chapter{Lebesgue測度}

\paragraph{本章の目標}
有界な窓で切ったときにLebesgue外測度 $\lambda^*$ とLebesgue内測度 $\lambda_*$ が一致する集合を
Lebesgue可測集合と定める．
そのうえで，Lebesgue可測集合全体が $\sigma$-加法族をなし，
Lebesgue測度が可算加法的な測度になることを確認する．
%最後に，内外正則性とCarathéodory条件による判定法を述べる．

\section{Lebesgue可測集合とLebesgue測度}

以下では，正数$R>0$と集合$A \subset \RR^d$に対し，
\[
Q_R:=[-R,R]^d, %\qquad(R>0)
\qquad A_R := A \cap Q_R
\]
と略記することがある．

\begin{definition}[Lebesgue可測集合とLebesgue測度]\label{def:lebesgue-measurable-local}
集合 $A \subset \RR^d$ がLebesgue可測であるとは
\[
\lambda_*(A\cap Q_R)=\lambda^*(A\cap Q_R)
\qquad (R>0)
\]
が成り立つことをいう．
Lebesgue可測集合全体を
\[
\calL_d:=\{A\subset\RR^d\mid A \text{ はLebesgue可測}\}
\]
と書く．
$A\in\calL_d$ のとき，
\[
\lambda(A):=\lambda^*(A)
\]
を $A$ のLebesgue測度という．
\end{definition}

% \begin{lemma}[有界集合の可測性]\label{lem:bounded-measurable-iff-inner-outer}
% 有界集合 $A\subset\RR^d$ について，次は同値である．
% \begin{enumerate}
%     \item $A$ はLebesgue可測である
%     \item $\lambda_*(A)=\lambda^*(A)$
% \end{enumerate}
% また，$A$ が有界ならば 
% %\[
% $\lambda_*(A)<\infty$
% %\]
% である．
% \end{lemma}
% \begin{proof}
% $A$ は有界なので，ある $R_0>0$ が存在して
% \[
% A\subset Q_{R_0}=[-R_0,R_0]^d
% \]
% となる．

% まず $A$ がLebesgue可測であるとする．
% 定義より
% \[
% \lambda_*(A\cap Q_{R_0})=\lambda^*(A\cap Q_{R_0})
% \]
% である．
% $A\cap Q_{R_0}=A$ だから
% \[
% \lambda_*(A)=\lambda^*(A)
% \]
% を得る．

% 逆に $\lambda_*(A)=\lambda^*(A)$ とする．
% 区間による外測度の分解（\cref{cor:interval-outer-measure-split}）を
% $A\subset Q_{R_0}$ に適用すると
% \[
% \lambda^*(Q_{R_0})
% =
% \lambda^*(A)+\lambda^*(Q_{R_0}\setminus A)
% \]
% である．

% 任意の $R>0$ に対し，
% \[
% A_R:=A\cap Q_R
% \]
% とおく．
% $A_R$ が内外一致することを示せばよい．
% $R>R_0$ ならば $A_R=A$ だから，仮定より内外一致する．
% 以下では $R\le R_0$ とする．

% コンパクト集合による外測度の分解
% （\cref{lem:compact-splits-outer-measure}）を
% コンパクト集合 $Q_R$ と集合 $A,\,Q_{R_0}\setminus A,\,Q_{R_0}$ にそれぞれ適用すると，
% \[
% \lambda^*(A)=\lambda^*(A_R)+\lambda^*(A\setminus Q_R),
% \]
% \[
% \lambda^*(Q_{R_0}\setminus A)
% =
% \lambda^*(Q_R\setminus A_R)
% +
% \lambda^*((Q_{R_0}\setminus A)\setminus Q_R),
% \]
% \[
% \lambda^*(Q_{R_0})
% =
% \lambda^*(Q_R)+\lambda^*(Q_{R_0}\setminus Q_R)
% \]
% である．ここで
% \[
% Q_{R_0}\setminus Q_R
% =
% (A\setminus Q_R)\cup((Q_{R_0}\setminus A)\setminus Q_R)
% \]
% だから，外測度の劣加法性より
% \[
% \lambda^*(Q_{R_0}\setminus Q_R)
% \le
% \lambda^*(A\setminus Q_R)
% +
% \lambda^*((Q_{R_0}\setminus A)\setminus Q_R)
% \]
% である．
% すべて $Q_{R_0}$ の部分集合なので外測度は有限であり，上の等式を比較すると
% \[
% \lambda^*(A_R)+\lambda^*(Q_R\setminus A_R)
% \le
% \lambda^*(Q_R)
% \]
% を得る．逆向きの不等式は外測度の劣加法性から従うので
% \[
% \lambda^*(Q_R)
% =
% \lambda^*(A_R)+\lambda^*(Q_R\setminus A_R)
% \]
% である．
% 最後に，第3章で示した区間の外測度の計算より $\lambda^*(Q_R)=|Q_R|$ であるから，
% 区間を用いた内測度表示を $A_R\subset Q_R$ に適用すれば
% \[
% \lambda_*(A_R)
% =
% \lambda^*(Q_R)-\lambda^*(Q_R\setminus A_R)
% =
% \lambda^*(A_R)
% \]
% を得る．
% よって $A_R=A\cap Q_R$ は内外一致する．
% $R>0$ は任意だから，$A$ はLebesgue可測である．

% 最後に，$A\subset Q_{R_0}$ と
% Lebesgue外測度の単調性（\cref{thm:lebesgue-outer-measure}）より
% \[
% \lambda_*(A)\le \lambda^*(A)\le \lambda^*(Q_{R_0})<\infty
% \]
% である．
% \end{proof}

% \begin{corollary}[有界集合の区間判定]\label{cor:bounded-measurable-interval-criterion}
% 有界集合 $A\subset\RR^d$ と，$A\subset Q$ を満たす有界区間 $Q$ に対して，
% 次は同値である．
% \begin{enumerate}
%     \item $A$ はLebesgue可測である
%     \item $\lambda^*(A)+\lambda^*(Q\setminus A)=|Q|$
% \end{enumerate}
% \end{corollary}
% \begin{proof}
% 有界集合の可測性（\cref{lem:bounded-measurable-iff-inner-outer}）と
% 区間による外測度の分解（\cref{cor:interval-outer-measure-split}）より，
% $A$ がLebesgue可測であることは
% \[
% \lambda^*(A)+\lambda^*(Q\setminus A)=|Q|
% \]
% と同値である．
% \end{proof}

\begin{theorem}[有限外測度集合の可測性]\label{thm:finite-outer-measurable-iff-inner-outer}
$\lambda^*(A)<\infty$ とする．このとき次は同値である．
\begin{enumerate}
    \item $A$ はLebesgue可測である
    \item $\lambda_*(A)=\lambda^*(A)$
\end{enumerate}
\end{theorem}
\begin{proof}
まず $A$ がLebesgue可測であるとする．
$\eps>0$ を任意に取る．
Lebesgue外測度の定義より，$A$ を覆う区間列 $\{I_n\}_{n=1}^{\infty}$ で
\[
\sum_{n=1}^{\infty}|I_n|<\lambda^*(A)+\eps
\]
となるものが取れる．
ある $N$ について
\[
\sum_{n>N}|I_n|<\eps
\]
である．さらに $R>0$ を十分大きく取って，
空でない $I_1,\dots,I_N$ がすべて $Q_R$ に含まれるようにする．
すると
\[
A\setminus Q_R\subset\bigcup_{n>N} I_n
\]
だから
\[
\lambda^*(A\setminus Q_R)<\eps.
\]
$A$ はLebesgue可測なので，$A\cap Q_R$ は内外一致する．
したがって，コンパクト集合 $K\subset A\cap Q_R$ で
\[
\lambda^*(K)>\lambda^*(A\cap Q_R)-\eps
\]
となるものが取れる．外測度の劣加法性より
\[
\lambda^*(A)
\le
\lambda^*(A\cap Q_R)+\lambda^*(A\setminus Q_R)
<
\lambda^*(K)+2\eps
\le
\lambda_*(A)+2\eps.
\]
$\eps>0$ は任意なので，Lebesgue内測度の基本性質
（\cref{prop:lebesgue-inner-basic}）と合わせて
\[
\lambda_*(A)=\lambda^*(A)
\]
である．

逆に $\lambda_*(A)=\lambda^*(A)$ とする．
任意の $R>0$ に対し，
\[
A_R:=A\cap Q_R
\]
が内外一致することを示す．
$\eps>0$ を任意に取る．内測度の定義より，コンパクト集合 $K\subset A$ で
\[
\lambda^*(K)>\lambda^*(A)-\eps
\]
となるものが取れる．
コンパクト集合による外測度の分解（\cref{lem:compact-splits-outer-measure}）を
コンパクト集合 $Q_R$ と集合 $A,\,K$ にそれぞれ適用すると，
\[
\lambda^*(A)=\lambda^*(A_R)+\lambda^*(A\setminus Q_R),
\]
\[
\lambda^*(K)
=
\lambda^*(K\cap Q_R)+\lambda^*(K\setminus Q_R)
\]
である．また $K\setminus Q_R\subset A\setminus Q_R$ だから
\[
\lambda^*(K\cap Q_R)
>
\lambda^*(A)-\lambda^*(A\setminus Q_R)-\eps
=
\lambda^*(A_R)-\eps.
\]
$K\cap Q_R$ は $A_R$ に含まれるコンパクト集合なので，
$\lambda_*(A_R)>\lambda^*(A_R)-\eps$ である．
$\eps>0$ は任意だから
\[
\lambda_*(A_R)\ge \lambda^*(A_R)
\]
を得る．逆向きの不等式はLebesgue内測度の基本性質であるから，
$A_R$ は内外一致する．
$R>0$ は任意だから，$A$ はLebesgue可測である．
\end{proof}

\begin{remark}
一般に$A$が有界集合ならば外測度有限 $\lambda^*(A)<\infty$ である．
しかし，逆
% \[
% \lambda^*(A)<\infty \implies A \text{ は有界}
% \]
は成り立たない．
例えば $\ZZ^d$ は非有界であるが，可算集合なので
%\[
$\lambda^*(\ZZ^d)=0$
%\]
である．

また，有限外測度集合の可測性
（\cref{thm:finite-outer-measurable-iff-inner-outer}）では
$\lambda^*(A)<\infty$ という仮定が本質的である．
後で示すLebesgue非可測集合 $V\subset[0,1]\subset\RR$ を用いて
\[
A:=V\cup[2,\infty)\subset\RR
\]
とおく．このとき任意の $N>2$ に対して $[2,N]\subset A$ だから
\[
\lambda_*(A)\ge \lambda^*([2,N])=N-2
\]
であり，$N\to\infty$ として $\lambda_*(A)=\infty$ である．
したがって $\lambda^*(A)=\infty$ でもあり，
\[
\lambda_*(A)=\lambda^*(A)=\infty
\]
が成り立つ．一方，
\[
A\cap[-1,1]=V
\]
であるから，$A$ はLebesgue可測ではない．
\end{remark}

% \begin{definition}[Lebesgue零集合]
% 集合 $N\subset\RR^d$ が
% \[
% \lambda^*(N)=0
% \]
% を満たすとき，$N$ をLebesgue零集合という．
% \end{definition}

\begin{theorem}[Lebesgue零集合]\label{thm:lebesgue-null-measurable}
Lebesgue零集合はLebesgue可測であり，Lebesgue測度は $0$ である．
\end{theorem}
\begin{proof}
任意の $R>0$ に対して，
$Q_R := [-R,R]^d$ と書く．
Lebesgue外測度の単調性
（\cref{thm:lebesgue-outer-measure}）より
\[
\lambda^*(N\cap Q_R)\le \lambda^*(N)=0
\]
である．
Lebesgue内測度の基本性質（\cref{prop:lebesgue-inner-basic}）より
\[
0\le \lambda_*(N\cap Q_R)
\le
\lambda^*(N\cap Q_R)
\]
である．
したがって
\[
\lambda_*(N\cap Q_R)
=
\lambda^*(N\cap Q_R)
=0
\]
だから，$N$ はLebesgue可測である．
また $\lambda(N)=\lambda^*(N)=0$ である．
\end{proof}

\begin{corollary}[可算集合]\label{cor:countable-lebesgue-null}
一点集合，可算集合，$\QQ^d$，$\ZZ^d$ はLebesgue可測であり，測度 $0$ をもつ．
%特に1次元では $\QQ$ と $\ZZ$ はLebesgue測度 $0$ である．
\end{corollary}
\begin{proof}
一点集合および可算集合のLebesgue外測度は $0$ である
（\cref{thm:countable-outer-measure-zero}）．
したがってLebesgue零集合の可測性（\cref{thm:lebesgue-null-measurable}）より従う．
\end{proof}

\begin{theorem}[Jordan可測集合]\label{thm:jordan-measurable-lebesgue}
Jordan可測集合 $A\subset\RR^d$ はLebesgue可測であり，
\[
\lambda(A)=m_J(A)
\]
が成り立つ．
\end{theorem}
\begin{proof}
    Jordan可測集合は有界なのでLebesgue外測度有限である（$\lambda^*(A) < \infty$）．よって$\lambda_*(A)=\lambda^*(A)$を示せばよい．
Jordan可測性より
\[
m_{J,*}(A)=m_J^*(A)=m_J(A)
\]
である．
第4章のJordan内測度との比較，Lebesgue内測度の基本性質，
第3章のJordan外測度との比較を順に用いると
\[
m_J(A)
=m_{J,*}(A)
\le \lambda_*(A)
\le \lambda^*(A)
\le m_J^*(A)
=m_J(A)
\]
となる．
したがってすべて等号であり，特に
\[
\lambda_*(A)=\lambda^*(A)=m_J(A)
\]
である．
よって $A$ はLebesgue可測であり，
\[
\lambda(A)=m_J(A)
\]
である．
\end{proof}
% \begin{proof}
% Jordan可測集合は有界なので，有界閉区間 $Q$ を
% \[
% A\subset Q
% \]
% となるように取る．
% $Q\setminus A$ もJordan可測であり，Jordan測度の有限加法性より
% \[
% m_J(A)+m_J(Q\setminus A)=m_J(Q)=|Q|
% \]
% である．
% 第3章で示したJordan測度とLebesgue外測度の一致より
% \[
% \lambda^*(A)+\lambda^*(Q\setminus A)=|Q|
% \]
% だから，有界集合の区間判定（\cref{cor:bounded-measurable-interval-criterion}）より
% $A$ はLebesgue可測である．
% さらに
% \[
% \lambda(A)=\lambda^*(A)=m_J(A)
% \]
% は，ふたたびJordan測度とLebesgue外測度の一致から従う．
% \end{proof}

\begin{example}
\[
\QQ\cap[0,1]
\]
は可算集合なのでLebesgue可測であり，
\[
\lambda_1(\QQ\cap[0,1])=0
\]
である．
一方，$[0,1]$ はJordan可測で $m_J([0,1])=1$ だから
\[
\lambda_1([0,1])=1
\]
である．
よって，後で示す完全加法性（\cref{thm:lebesgue-measure-countable-additivity}）から従う有限加法性より
\[
\lambda_1([0,1]\setminus\QQ)=1
\]
となる．
\end{example}

\section{Lebesgue可測集合の性質}

\begin{definition}[$\sigma$-加法族]
集合 $X$ の部分集合族 $\calS\subset\calP(X)$ が
$X$ 上の $\sigma$-加法族であるとは，次の3条件を満たすことをいう．
\begin{enumerate}
    \item $X\in\calS$
    \item $A\in\calS$ ならば $X\setminus A\in\calS$
    \item $\{A_n\}_{n=1}^{\infty}\subset\calS$ ならば
    \[
    \bigcup_{n=1}^{\infty}A_n\in\calS
    \]
\end{enumerate}
\end{definition}

\begin{lemma}[有限外測度集合の近似判定法]\label{lem:finite-outer-compact-criterion}
$\lambda^*(A)<\infty$ とする．このとき次は同値である．
\begin{enumerate}
    \item $A$ はLebesgue可測である
    \item 任意の $\eps>0$ に対して，コンパクト集合 $K\subset A$ で
    \[
    \lambda^*(A\setminus K)<\eps
    \]
    となるものが取れる
    \item 任意の $\eps>0$ に対して，開集合 $U\supset A$ で
    \[
    \lambda^*(U\setminus A)<\eps
    \]
    となるものが取れる
\end{enumerate}
\end{lemma}
\begin{proof}
まず (1) と (2) が同値であることを示す．
$A$ がLebesgue可測であるとする．有限外測度集合の可測性
（\cref{thm:finite-outer-measurable-iff-inner-outer}）より
\[
\lambda_*(A)=\lambda^*(A)
\]
である．内測度の定義から，任意の $\eps>0$ に対して
コンパクト集合 $K\subset A$ で
\[
\lambda^*(K)>\lambda^*(A)-\eps
\]
となるものが取れる．
コンパクト集合による外測度の分解
（\cref{lem:compact-splits-outer-measure}）を $B=A$ と $C=K$ に適用すると
\[
\lambda^*(A)=\lambda^*(K)+\lambda^*(A\setminus K)
\]
である．したがって
\[
\lambda^*(A\setminus K)<\eps
\]
である．

逆に，任意の $\eps>0$ に対して上のようなコンパクト集合 $K\subset A$ が取れるとする．
コンパクト集合による外測度の分解
（\cref{lem:compact-splits-outer-measure}）より
\[
\lambda^*(A)=\lambda^*(K)+\lambda^*(A\setminus K)
\]
である．また $K\subset A$ はコンパクトだから，内測度の定義より
\[
\lambda^*(K)\le \lambda_*(A)
\]
である．よって
\[
\lambda^*(A)<\lambda_*(A)+\eps
\]
を得る．$\eps>0$ は任意であり，Lebesgue内測度の基本性質
（\cref{prop:lebesgue-inner-basic}）より $\lambda_*(A)\le\lambda^*(A)$ だから
\[
\lambda_*(A)=\lambda^*(A)
\]
である．有限外測度集合の可測性
（\cref{thm:finite-outer-measurable-iff-inner-outer}）より，$A$ はLebesgue可測である．

次に (1) と (3) が同値であることを示す．
まず $A$ がLebesgue可測であるとする．
有限外測度集合の可測性
（\cref{thm:finite-outer-measurable-iff-inner-outer}）より
\[
\lambda_*(A)=\lambda^*(A)
\]
である．$\eps>0$ を任意に取る．
開集合による外正則性
（\cref{thm:lebesgue-outer-regularity-open}）より，
$A\subset U$ を満たす開集合 $U$ で
\[
\lambda^*(U)<\lambda^*(A)+\frac{\eps}{2}
\]
となるものが取れる．また内測度の定義より，
コンパクト集合 $K\subset A$ で
\[
\lambda^*(K)>\lambda^*(A)-\frac{\eps}{2}
\]
となるものが取れる．
$K\subset U$ であり，$\lambda^*(U)<\infty$ だから，
コンパクト集合による外測度の分解
（\cref{lem:compact-splits-outer-measure}）を $B=U$ と $C=K$ に適用すると
\[
\lambda^*(U)=\lambda^*(K)+\lambda^*(U\setminus K)
\]
である．$U\setminus A\subset U\setminus K$ なので
\[
\lambda^*(U\setminus A)
\le
\lambda^*(U\setminus K)
=
\lambda^*(U)-\lambda^*(K)
<
\eps
\]
である．したがって (3) が成り立つ．

逆に (3) を仮定する．$\eps>0$ を任意に取る．
仮定より，$A\subset U$ を満たす開集合 $U$ で
\[
\lambda^*(U\setminus A)<\frac{\eps}{3}
\]
となるものが取れる．このとき
\[
\lambda^*(U)\le \lambda^*(A)+\lambda^*(U\setminus A)<\infty
\]
である．開集合による外正則性
（\cref{thm:lebesgue-outer-regularity-open}）より，
$U\setminus A\subset V$ を満たす開集合 $V$ で
\[
\lambda^*(V)<\lambda^*(U\setminus A)+\frac{\eps}{3}
<\frac{2\eps}{3}
\]
となるものが取れる．
また，有限外測度の開集合をコンパクト集合で内側から近似できること
（\cref{lem:finite-open-inner-compact}）より，コンパクト集合 $L\subset U$ で
\[
\lambda^*(L)>\lambda^*(U)-\frac{\eps}{3}
\]
となるものが取れる．
有限外測度の開集合とコンパクト集合の分解
（\cref{lem:finite-open-compact-split}）を $U$ と $L$ に適用すると
\[
\lambda^*(U)=\lambda^*(L)+\lambda^*(U\setminus L)
\]
だから
\[
\lambda^*(U\setminus L)<\frac{\eps}{3}
\]
である．

\[
K:=L\setminus V
\]
とおく．$L$ はコンパクトで $V$ は開集合だから，$K$ はコンパクトである．
さらに $L\subset U$ かつ $V\supset U\setminus A$ なので
\[
K\subset U\setminus V\subset A
\]
である．また
\[
U\subset K\cup V\cup(U\setminus L)
\]
だから，外測度の劣加法性
（\cref{thm:lebesgue-outer-measure}）より
\[
\lambda^*(A)
\le
\lambda^*(U)
\le
\lambda^*(K)+\lambda^*(V)+\lambda^*(U\setminus L)
<
\lambda^*(K)+\eps.
\]
したがって
\[
\lambda_*(A)\ge \lambda^*(K)>\lambda^*(A)-\eps
\]
である．$\eps>0$ は任意なので，
Lebesgue内測度の基本性質
（\cref{prop:lebesgue-inner-basic}）と合わせて
\[
\lambda_*(A)=\lambda^*(A)
\]
を得る．有限外測度集合の可測性
（\cref{thm:finite-outer-measurable-iff-inner-outer}）より，$A$ はLebesgue可測である．
\end{proof}

\subsection{有界閉区間上のLebesgue可測集合}

有界な窓の中では，Lebesgue可測性は単に内外測度の一致として扱える．
以下，$Q\subset\RR^d$ を有界閉区間とし，
\[
\calL(Q):=\{A\subset Q\mid \lambda_*(A)=\lambda^*(A)\}
\]
と書く．

\begin{lemma}[有界閉区間内の補集合]\label{lem:local-measurable-complement}
$A\in\calL(Q)$ ならば $Q\setminus A\in\calL(Q)$ である．
\end{lemma}
\begin{proof}
$A\in\calL(Q)$ だから
\[
\lambda_*(A)=\lambda^*(A)
\]
である．区間による外測度の分解
（\cref{cor:interval-outer-measure-split}）より
\[
|Q|=\lambda^*(A)+\lambda^*(Q\setminus A)
\]
である．同じ区間による外測度の分解
（\cref{cor:interval-outer-measure-split}）を $Q\setminus A\subset Q$ に適用すると，
この等式は
\[
\lambda_*(Q\setminus A)=\lambda^*(Q\setminus A)
\]
と同値である．したがって $Q\setminus A\in\calL(Q)$ である．
\end{proof}

\begin{lemma}[有界閉区間内の有限和と有限加法性]\label{lem:local-measurable-finite-union-additivity}
$A,B\in\calL(Q)$ とする．このとき $A\cup B\in\calL(Q)$ である．
さらに $A\cap B=\emptyset$ ならば
\[
\lambda^*(A\cup B)=\lambda^*(A)+\lambda^*(B)
\]
である．
\end{lemma}
\begin{proof}
まず有限和に関する閉性を示す．
$A,B\subset Q$ なので，$\lambda^*(A),\lambda^*(B)<\infty$ である．
$A,B\in\calL(Q)$ と有限外測度集合の可測性
（\cref{thm:finite-outer-measurable-iff-inner-outer}）より，
$A,B$ はLebesgue可測である．
有限外測度集合の近似判定法
（\cref{lem:finite-outer-compact-criterion}）より，任意の $\eps>0$ に対して
コンパクト集合 $K_A\subset A,\ K_B\subset B$ で
\[
\lambda^*(A\setminus K_A)<\frac{\eps}{2},
\qquad
\lambda^*(B\setminus K_B)<\frac{\eps}{2}
\]
となるものが取れる．
このとき
\[
K:=K_A\cup K_B
\]
は $A\cup B$ に含まれるコンパクト集合であり，
\[
(A\cup B)\setminus K
\subset
(A\setminus K_A)\cup(B\setminus K_B)
\]
である．Lebesgue外測度の性質
（\cref{thm:lebesgue-outer-measure}）より
\[
\lambda^*((A\cup B)\setminus K)<\eps
\]
である．したがって有限外測度集合の近似判定法
（\cref{lem:finite-outer-compact-criterion}）より $A\cup B$ はLebesgue可測である．
また $A\cup B\subset Q$ だから，有限外測度集合の可測性
（\cref{thm:finite-outer-measurable-iff-inner-outer}）より
\[
\lambda_*(A\cup B)=\lambda^*(A\cup B)
\]
である．よって $A\cup B\in\calL(Q)$ である．

次に $A\cap B=\emptyset$ とする．
外測度の劣加法性（\cref{thm:lebesgue-outer-measure}）より
\[
\lambda^*(A\cup B)\le \lambda^*(A)+\lambda^*(B)
\]
である．逆向きの不等式を示す．$\eps>0$ を任意に取る．
$A,B\in\calL(Q)$ だから，内測度の定義よりコンパクト集合
$K_A\subset A,\ K_B\subset B$ で
\[
\lambda^*(K_A)>\lambda^*(A)-\frac{\eps}{2},
\qquad
\lambda^*(K_B)>\lambda^*(B)-\frac{\eps}{2}
\]
となるものが取れる．
$A\cap B=\emptyset$ なので $K_A$ と $K_B$ は交わらない．
空でない場合には2つのコンパクト集合は正の距離だけ離れているから，
正距離集合に対する外測度の加法性
（\cref{lem:positive-distance-additive}）より
\[
\lambda^*(K_A\cup K_B)
=
\lambda^*(K_A)+\lambda^*(K_B)
\]
である．片方が空の場合もこの等式は自明である．
したがって
\[
\lambda^*(A\cup B)
\ge
\lambda^*(K_A\cup K_B)
>
\lambda^*(A)+\lambda^*(B)-\eps
\]
である．$\eps>0$ は任意だから
\[
\lambda^*(A\cup B)\ge \lambda^*(A)+\lambda^*(B)
\]
を得る．
\end{proof}

\begin{corollary}[有界閉区間内の集合演算]\label{cor:local-measurable-set-operations}
$A,B\in\calL(Q)$ ならば
\[
A\cap B,\qquad A\setminus B
\]
も $\calL(Q)$ に属する．
\end{corollary}
\begin{proof}
有界閉区間内の補集合
（\cref{lem:local-measurable-complement}）より $Q\setminus A,Q\setminus B\in\calL(Q)$ である．
有界閉区間内の有限和に関する閉性
（\cref{lem:local-measurable-finite-union-additivity}）より
\[
(Q\setminus A)\cup(Q\setminus B)\in\calL(Q),
\qquad
(Q\setminus A)\cup B\in\calL(Q)
\]
である．もう一度補集合に関する閉性を用いると
\[
A\cap B
=
Q\setminus\bigl((Q\setminus A)\cup(Q\setminus B)\bigr)\in\calL(Q)
\]
および
\[
A\setminus B
=
Q\setminus\bigl((Q\setminus A)\cup B\bigr)\in\calL(Q)
\]
を得る．
\end{proof}

\begin{theorem}[有界閉区間内の $\sigma$-加法族性]\label{thm:local-measurable-sigma-algebra}
$\calL(Q)$ は $Q$ 上の $\sigma$-加法族である．
さらに，互いに素な列 $\{A_n\}_{n=1}^{\infty}\subset\calL(Q)$ に対して
\[
\lambda^*\left(\bigsqcup_{n=1}^{\infty}A_n\right)
=
\sum_{n=1}^{\infty}\lambda^*(A_n)
\]
が成り立つ．
\end{theorem}
\begin{proof}
まず $Q\in\calL(Q)$ を確認する．
$Q$ はコンパクト集合なので，コンパクト集合では内外が一致する
（\cref{prop:compact-inner-outer-equal}）より
\[
\lambda_*(Q)=\lambda^*(Q)
\]
である．また有界閉区間内の補集合
（\cref{lem:local-measurable-complement}）より $\emptyset=Q\setminus Q\in\calL(Q)$ である．

補集合に関する閉性は有界閉区間内の補集合
（\cref{lem:local-measurable-complement}）で示した．
したがって，あとは可算和に関する閉性を示せばよい．

まず $\{A_n\}_{n=1}^{\infty}\subset\calL(Q)$ が互いに素である場合を考える．
\[
A:=\bigcup_{n=1}^{\infty}A_n,
\qquad
S_N:=\bigcup_{n=1}^{N}A_n
\]
とおく．有界閉区間内の有限和に関する閉性と有限加法性
（\cref{lem:local-measurable-finite-union-additivity}）を繰り返し用いると
\[
S_N\in\calL(Q),
\qquad
\lambda^*(S_N)=\sum_{n=1}^{N}\lambda^*(A_n)
\]
である．単調性より
\[
\sum_{n=1}^{N}\lambda^*(A_n)=\lambda^*(S_N)\le\lambda^*(A)
\]
であるから，$N\to\infty$ として
\[
\sum_{n=1}^{\infty}\lambda^*(A_n)\le\lambda^*(A)
\]
を得る．逆向きの不等式は外測度の可算劣加法性
（\cref{thm:lebesgue-outer-measure}）から従う．よって
\[
\lambda^*(A)=\sum_{n=1}^{\infty}\lambda^*(A_n)
\]
である．

$A\subset Q$ なので $\lambda^*(A)<\infty$ である．
任意の $\eps>0$ を取る．上の等式より，ある $N$ が存在して
\[
\lambda^*(S_N)>\lambda^*(A)-\frac{\eps}{2}
\]
である．$S_N\in\calL(Q)$ かつ $S_N\subset Q$ だから，有限外測度集合の可測性
（\cref{thm:finite-outer-measurable-iff-inner-outer}）より $S_N$ はLebesgue可測である．
したがって有限外測度集合の近似判定法
（\cref{lem:finite-outer-compact-criterion}）より，
コンパクト集合 $K\subset S_N$ で
\[
\lambda^*(S_N\setminus K)<\frac{\eps}{2}
\]
となるものが取れる．コンパクト集合による外測度の分解
（\cref{lem:compact-splits-outer-measure}）より
\[
\lambda^*(S_N)=\lambda^*(K)+\lambda^*(S_N\setminus K)
\]
だから
\[
\lambda^*(K)>\lambda^*(S_N)-\frac{\eps}{2}
>\lambda^*(A)-\eps
\]
である．
したがって有限外測度集合の近似判定法
（\cref{lem:finite-outer-compact-criterion}）より $A$ はLebesgue可測である．
さらに $A\subset Q$ なので，有限外測度集合の可測性
（\cref{thm:finite-outer-measurable-iff-inner-outer}）より $A\in\calL(Q)$ である．
互いに素な場合の可算和に関する閉性と加法性が示された．

一般の列 $\{A_n\}_{n=1}^{\infty}\subset\calL(Q)$ に対しては
\[
B_1:=A_1,\qquad
B_n:=A_n\setminus\bigcup_{j=1}^{n-1}A_j
\quad(n\ge2)
\]
とおく．有界閉区間内の有限和に関する閉性
（\cref{lem:local-measurable-finite-union-additivity}）と
有界閉区間内の集合演算
（\cref{cor:local-measurable-set-operations}）より，すべての $n$ について
$B_n\in\calL(Q)$ である．また $\{B_n\}$ は互いに素であり，
\[
\bigcup_{n=1}^{\infty}B_n=\bigcup_{n=1}^{\infty}A_n
\]
である．したがって上で示した互いに素な場合を適用すれば，
\[
\bigcup_{n=1}^{\infty}A_n\in\calL(Q)
\]
を得る．
\end{proof}

\subsection{全空間上のLebesgue可測集合}

\begin{remark}
    次の性質を使用して，Lebesgue可測集合の $\sigma$-加法族性を示す．
    \begin{itemize}
        \item $\RR^d$ の任意の有界閉区間 $Q$ に対して，$\calL(Q)$ は $Q$ 上の $\sigma$-加法族である．
        \item $A\subset\RR^d$ がLebesgue可測であるとは，任意の有界閉区間 $Q$ に対して $A\cap Q\in\calL(Q)$ であることと同値である．
    \end{itemize}
\end{remark}

\begin{theorem}[Lebesgue可測集合の $\sigma$-加法族性]\label{thm:lebesgue-measurable-sigma-algebra}
Lebesgue可測集合全体 $\calL_d$ は $\RR^d$ 上の $\sigma$-加法族である．
\end{theorem}
\begin{proof}
まず $\RR^d\in\calL_d$ である．実際，任意の $R>0$ に対して
\[
\RR^d\cap Q_R=Q_R
\]
であり，$Q_R$ はコンパクト集合なので，コンパクト集合では内外が一致する
（\cref{prop:compact-inner-outer-equal}）より内外一致する．
また $\emptyset\in\calL_d$ も同様である．

次に $A\in\calL_d$ とする．任意の $R>0$ に対して，
$A\cap Q_R\in\calL(Q_R)$ である．有界閉区間内の $\sigma$-加法族性
（\cref{thm:local-measurable-sigma-algebra}）より
\[
Q_R\setminus(A\cap Q_R)\in\calL(Q_R)
\]
である．これは
\[
(\RR^d\setminus A)\cap Q_R
=
Q_R\setminus(A\cap Q_R)
\]
に等しい．したがって任意の $R>0$ で
\[
\lambda_*\bigl((\RR^d\setminus A)\cap Q_R\bigr)
=
\lambda^*\bigl((\RR^d\setminus A)\cap Q_R\bigr)
\]
が成り立つので，$\RR^d\setminus A\in\calL_d$ である．

最後に $\{A_n\}_{n=1}^{\infty}\subset\calL_d$ とする．
任意の $R>0$ に対して，各 $A_n\cap Q_R$ は $\calL(Q_R)$ に属する．
有界閉区間内の $\sigma$-加法族性
（\cref{thm:local-measurable-sigma-algebra}）より
\[
\bigcup_{n=1}^{\infty}(A_n\cap Q_R)\in\calL(Q_R)
\]
である．ところが
\[
\left(\bigcup_{n=1}^{\infty}A_n\right)\cap Q_R
=
\bigcup_{n=1}^{\infty}(A_n\cap Q_R)
\]
である．したがって任意の $R>0$ で内外一致が成り立つので
\[
\bigcup_{n=1}^{\infty}A_n\in\calL_d
\]
である．
\end{proof}

\begin{corollary}[集合演算]\label{cor:lebesgue-measurable-set-operations}
$A,B\in\calL_d$ ならば
\[
A\cup B,\qquad A\cap B,\qquad A\setminus B
\]
も $\calL_d$ に属する．また $\{A_n\}_{n=1}^{\infty}\subset\calL_d$ ならば
\[
\bigcap_{n=1}^{\infty}A_n\in\calL_d
\]
である．
\end{corollary}
\begin{proof}
Lebesgue可測集合の $\sigma$-加法族性
（\cref{thm:lebesgue-measurable-sigma-algebra}）から，補集合と可算和に関する閉性が従う．
あとはDe Morganの法則と
\[
A\setminus B=A\cap(\RR^d\setminus B)
\]
を用いればよい．
\end{proof}

\begin{theorem}[開集合と閉集合]\label{thm:open-closed-lebesgue-measurable}
$\RR^d$ の開集合と閉集合はLebesgue可測である．
\end{theorem}
\begin{proof}
まず閉集合 $F\subset\RR^d$ を取る．任意の $R>0$ に対して
\[
F\cap Q_R
\]
はコンパクト集合である．コンパクト集合では内外が一致する
（\cref{prop:compact-inner-outer-equal}）より
\[
\lambda_*(F\cap Q_R)=\lambda^*(F\cap Q_R)
\]
である．したがって $F$ はLebesgue可測である．

開集合 $G$ については，$\RR^d\setminus G$ が閉集合なのでLebesgue可測である．
Lebesgue可測集合の $\sigma$-加法族性
（\cref{thm:lebesgue-measurable-sigma-algebra}）より，その補集合である $G$ もLebesgue可測である．
\end{proof}

\section{Lebesgue測度の性質} \label{sec:lebesgue-measure-properties}

\begin{definition}[完全加法的測度]\label{def:countably-additive-measure}
集合 $X$ と $X$ 上の $\sigma$-加法族 $\calS$ に対し，写像
\[
\mu:\calS\to[0,\infty]
\]
が $(X,\calS)$ 上の（完全加法的）測度であるとは，次の2条件を満たすことをいう．
\begin{enumerate}
    \item $\mu(\emptyset)=0$
    \item 互いに素な列 $\{A_n\}_{n=1}^{\infty}\subset\calS$ に対して
    \[
    \mu\left(\bigsqcup_{n=1}^{\infty}A_n\right)
    =
    \sum_{n=1}^{\infty}\mu(A_n)
    \]
\end{enumerate}
%第2条件を完全加法性という．
\end{definition}

\begin{proposition}[非負性と空集合]\label{prop:lebesgue-measure-nonnegative}
任意の $A\in\calL_d$ に対して
\[
0\le\lambda(A)\le\infty
\]
である．また
\[
\lambda(\emptyset)=0
\]
である．
\end{proposition}
\begin{proof}
Lebesgue測度の定義より $\lambda(A)=\lambda^*(A)$ である．
Lebesgue外測度の非負性
（\cref{thm:lebesgue-outer-measure}）より，$\lambda^*(A)\in[0,\infty]$ である．
また同じ外測度の基本性質より
\[
\lambda(\emptyset)=\lambda^*(\emptyset)=0
\]
である．
\end{proof}

\begin{lemma}[窓による測度の表示]\label{lem:lebesgue-measure-window}
任意の $A\in\calL_d$ について
\[
\lambda(A)=\sup_{m\ge1}\lambda(A\cap Q_m).
\]
\end{lemma}
\begin{proof}
右辺が左辺以下であることはLebesgue外測度の単調性から従う．
逆向きについて，右辺を $M$ とおく．$M=\infty$ のときはただちに従うので，
$M<\infty$ とする．
\[
E_1:=A\cap Q_1,\qquad
E_m:=A\cap(Q_m\setminus Q_{m-1})\quad(m\ge2)
\]
とおくと，$E_m$ は互いに素なLebesgue可測集合であり，
$A=\bigcup_{m=1}^{\infty}E_m$ である．
有界閉区間内の有限加法性
（\cref{lem:local-measurable-finite-union-additivity}）を
$Q_N$ の中の互いに素な集合 $E_1,\dots,E_N\in\calL(Q_N)$ に適用すると
\[
\lambda(A\cap Q_N)=\sum_{m=1}^{N}\lambda(E_m)
\]
である．したがって
\[
\sum_{m=1}^{\infty}\lambda(E_m)=M
\]
である．外測度の可算劣加法性（\cref{thm:lebesgue-outer-measure}）より
\[
\lambda(A)=\lambda^*(A)\le
\sum_{m=1}^{\infty}\lambda^*(E_m)
=
\sum_{m=1}^{\infty}\lambda(E_m)=M
\]
となり，主張が従う．
\end{proof}

\begin{lemma}[有限加法性]\label{lem:lebesgue-measure-finite-additivity}
$A,B\in\calL_d$ が互いに素ならば
\[
\lambda(A\cup B)=\lambda(A)+\lambda(B)
\]
である．
\end{lemma}
\begin{proof}
任意の $m\ge1$ に対して，
$A\cap Q_m$ と $B\cap Q_m$ は $\calL(Q_m)$ に属し，互いに素である．
有界閉区間内の有限加法性
（\cref{lem:local-measurable-finite-union-additivity}）より
\[
\lambda\bigl((A\cup B)\cap Q_m\bigr)
=
\lambda(A\cap Q_m)+\lambda(B\cap Q_m)
\]
である．$m$ に関して上限を取ると，
\cref{lem:lebesgue-measure-window} と単調列の上限の性質から
\[
\lambda(A\cup B)=\lambda(A)+\lambda(B)
\]
を得る．
\end{proof}

\begin{theorem}[完全加法性]\label{thm:lebesgue-measure-countable-additivity}
互いに素な列 $\{A_n\}_{n=1}^{\infty}\subset\calL_d$ に対して
\[
\lambda\left(\bigsqcup_{n=1}^{\infty}A_n\right)
=
\sum_{n=1}^{\infty}\lambda(A_n)
\]
が成り立つ．
\end{theorem}
\begin{proof}
Lebesgue可測集合の $\sigma$-加法族性
（\cref{thm:lebesgue-measurable-sigma-algebra}）より
\[
A:=\bigsqcup_{n=1}^{\infty}A_n
\]
はLebesgue可測である．
任意の $m\ge1$ に対して，
$\{A_n\cap Q_m\}_{n=1}^{\infty}\subset\calL(Q_m)$ は互いに素である．
有界閉区間内の完全加法性
（\cref{thm:local-measurable-sigma-algebra}）より
\[
\lambda(A\cap Q_m)
=
\sum_{n=1}^{\infty}\lambda(A_n\cap Q_m)
\le
\sum_{n=1}^{\infty}\lambda(A_n)
\]
である．$m$ に関して上限を取り，
\cref{lem:lebesgue-measure-window} を用いれば
\[
\lambda(A)\le\sum_{n=1}^{\infty}\lambda(A_n)
\]
である．

逆向きの不等式を示す．任意の $N$ に対して，
\cref{lem:lebesgue-measure-window} と単調列の上限の性質より
\[
\sum_{n=1}^{N}\lambda(A_n)
=
\sup_{m\ge1}\sum_{n=1}^{N}\lambda(A_n\cap Q_m)
\]
である．したがって，有界閉区間内の完全加法性
（\cref{thm:local-measurable-sigma-algebra}）より
\[
\sum_{n=1}^{N}\lambda(A_n)
\le
\sup_{m\ge1}\sum_{n=1}^{\infty}\lambda(A_n\cap Q_m)
=
\sup_{m\ge1}\lambda(A\cap Q_m)
=
\lambda(A)
\]
である．$N\to\infty$ とすれば
\[
\sum_{n=1}^{\infty}\lambda(A_n)\le\lambda(A)
\]
を得る．
\end{proof}

\begin{theorem}[単調性と劣加法性]\label{thm:lebesgue-measure-basic}
Lebesgue測度は次を満たす．
\begin{enumerate}
    \item $A,B\in\calL_d$ かつ $A\subset B$ ならば
    \[
    \lambda(A)\le\lambda(B)
    \]
    \item $\{A_n\}_{n=1}^{\infty}\subset\calL_d$ ならば
    \[
    \lambda\left(\bigcup_{n=1}^{\infty}A_n\right)
    \le
    \sum_{n=1}^{\infty}\lambda(A_n)
    \]
\end{enumerate}
\end{theorem}
\begin{proof}
\begin{enumerate}
    \item Lebesgue測度の定義と外測度の単調性
    （\cref{thm:lebesgue-outer-measure}）より
    \[
    \lambda(A)=\lambda^*(A)\le\lambda^*(B)=\lambda(B)
    \]
    である．

    \item Lebesgue可測集合の $\sigma$-加法族性
    （\cref{thm:lebesgue-measurable-sigma-algebra}）より
    $\bigcup_{n=1}^{\infty}A_n$ はLebesgue可測である．
    したがって，外測度の可算劣加法性
    （\cref{thm:lebesgue-outer-measure}）より
    \[
    \lambda\left(\bigcup_{n=1}^{\infty}A_n\right)
    =
    \lambda^*\left(\bigcup_{n=1}^{\infty}A_n\right)
    \le
    \sum_{n=1}^{\infty}\lambda^*(A_n)
    =
    \sum_{n=1}^{\infty}\lambda(A_n)
    \]
    である．
\end{enumerate}
\end{proof}

\begin{theorem}[下からの連続性]\label{thm:lebesgue-measure-continuity-from-below}
$\{A_n\}_{n=1}^{\infty}\subset\calL_d$ が
\[
A_1\subset A_2\subset \cdots
\]
を満たすとし，
\[
A:=\bigcup_{n=1}^{\infty}A_n
\]
とおく．このとき
\[
\lambda(A)=\lim_{n\to\infty}\lambda(A_n)
\]
である．
\end{theorem}
\begin{proof}
$A_0:=\emptyset$ とおき，
\[
B_n:=A_n\setminus A_{n-1}\qquad(n\ge1)
\]
と定める．Lebesgue可測集合の集合演算
（\cref{cor:lebesgue-measurable-set-operations}）より $B_n\in\calL_d$ であり，
$\{B_n\}$ は互いに素である．また
\[
A_N=\bigsqcup_{n=1}^{N}B_n,
\qquad
A=\bigsqcup_{n=1}^{\infty}B_n
\]
である．完全加法性（\cref{thm:lebesgue-measure-countable-additivity}）と
有限加法性（\cref{lem:lebesgue-measure-finite-additivity}）より
\[
\lambda(A)=\sum_{n=1}^{\infty}\lambda(B_n),
\qquad
\lambda(A_N)=\sum_{n=1}^{N}\lambda(B_n)
\]
である．したがって
\[
\lambda(A)=\lim_{N\to\infty}\lambda(A_N)
\]
を得る．
\end{proof}

\begin{theorem}[上からの連続性]\label{thm:lebesgue-measure-continuity-from-above}
$\{A_n\}_{n=1}^{\infty}\subset\calL_d$ が
\[
A_1\supset A_2\supset \cdots
\]
を満たし，$\lambda(A_1)<\infty$ とする．
\[
A:=\bigcap_{n=1}^{\infty}A_n
\]
とおくと
\[
\lambda(A)=\lim_{n\to\infty}\lambda(A_n)
\]
である．
\end{theorem}
\begin{proof}
\[
C_n:=A_1\setminus A_n
\]
とおく．Lebesgue可測集合の集合演算
（\cref{cor:lebesgue-measurable-set-operations}）より $C_n\in\calL_d$ であり，
\[
C_1\subset C_2\subset\cdots
\]
である．また
\[
\bigcup_{n=1}^{\infty}C_n=A_1\setminus A
\]
である．下からの連続性
（\cref{thm:lebesgue-measure-continuity-from-below}）より
\[
\lambda(A_1\setminus A)
=
\lim_{n\to\infty}\lambda(C_n)
\]
である．

有限加法性（\cref{lem:lebesgue-measure-finite-additivity}）より
\[
\lambda(A_1)=\lambda(A_n)+\lambda(C_n),
\qquad
\lambda(A_1)=\lambda(A)+\lambda(A_1\setminus A)
\]
である．仮定より $\lambda(A_1)<\infty$ なので，ここでは通常の実数の差として
\[
\lambda(A_n)=\lambda(A_1)-\lambda(C_n),
\qquad
\lambda(A)=\lambda(A_1)-\lambda(A_1\setminus A)
\]
と書ける．したがって
\[
\lim_{n\to\infty}\lambda(A_n)=\lambda(A)
\]
である．
\end{proof}

\begin{proposition}[平行移動不変性]\label{prop:lebesgue-measure-translation-invariant}
$A\in\calL_d$ と $c\in\RR^d$ に対して
\[
A+c:=\{x+c\mid x\in A\}
\]
もLebesgue可測であり，
\[
\lambda(A+c)=\lambda(A)
\]
である．
\end{proposition}
\begin{proof}
まず内測度も平行移動で不変であることを確認する．
コンパクト集合 $K\subset A$ とコンパクト集合 $K+c\subset A+c$ は一対一に対応する．
Lebesgue外測度の平行移動不変性
（\cref{prop:lebesgue-outer-translation-invariant}）より
\[
\lambda^*(K+c)=\lambda^*(K)
\]
であるから
\[
\lambda_*(A+c)=\lambda_*(A)
\]
である．

次に $A+c$ がLebesgue可測であることを示す．
任意の $R>0$ を取る．
\[
B:=A\cap(Q_R-c)
\]
とおく．$Q_R-c$ は有界閉区間なので，開集合と閉集合の可測性
（\cref{thm:open-closed-lebesgue-measurable}）よりLebesgue可測である．
Lebesgue可測集合の集合演算
（\cref{cor:lebesgue-measurable-set-operations}）より $B$ もLebesgue可測である．
また $B$ は有界なので $\lambda^*(B)<\infty$ である．有限外測度集合の可測性
（\cref{thm:finite-outer-measurable-iff-inner-outer}）より
\[
\lambda_*(B)=\lambda^*(B)
\]
である．上で示した内測度の平行移動不変性と
Lebesgue外測度の平行移動不変性
（\cref{prop:lebesgue-outer-translation-invariant}）より
\[
\lambda_*(B+c)=\lambda^*(B+c)
\]
である．ところが
\[
B+c=(A+c)\cap Q_R
\]
だから，任意の $R>0$ で内外一致が成り立つ．
よって $A+c$ はLebesgue可測である．

最後に，Lebesgue外測度の平行移動不変性
（\cref{prop:lebesgue-outer-translation-invariant}）より
\[
\lambda(A+c)=\lambda^*(A+c)=\lambda^*(A)=\lambda(A)
\]
である．
\end{proof}

\begin{theorem}[外正則性]\label{thm:lebesgue-measure-outer-regularity}
$A\in\calL_d$ に対して
\[
\lambda(A)
=
\inf\{\lambda(G)\mid A\subset G,\ G\text{ は開集合}\}
\]
が成り立つ．
\end{theorem}
\begin{proof}
開集合はLebesgue可測である
（\cref{thm:open-closed-lebesgue-measurable}）から，開集合 $G$ について
\[
\lambda(G)=\lambda^*(G)
\]
である．また $\lambda(A)=\lambda^*(A)$ である．
したがって，開集合による外正則性
（\cref{thm:lebesgue-outer-regularity-open}）より
\[
\lambda(A)
=
\lambda^*(A)
=
\inf\{\lambda^*(G)\mid A\subset G,\ G\text{ は開集合}\}
=
\inf\{\lambda(G)\mid A\subset G,\ G\text{ は開集合}\}
\]
を得る．
\end{proof}

\begin{theorem}[内正則性]\label{thm:lebesgue-measure-inner-regularity}
$A\in\calL_d$ に対して
\[
\lambda(A)
=
\sup\{\lambda(K)\mid K\subset A,\ K\text{ はコンパクト}\}
\]
が成り立つ．
\end{theorem}
\begin{proof}
コンパクト集合はLebesgue可測であり，そのLebesgue測度は外測度に等しいので，
右辺は内測度の定義そのものから $\lambda_*(A)$ である．
したがって，$\lambda(A)<\infty$ のときは，有限外測度集合の可測性
（\cref{thm:finite-outer-measurable-iff-inner-outer}）より
\[
\lambda(A)=\lambda^*(A)=\lambda_*(A)
\]
となり，結論が従う．

次に $\lambda(A)=\infty$ とする．
窓による測度の表示（\cref{lem:lebesgue-measure-window}）より
\[
\sup_{n\ge1}\lambda(A\cap Q_n)=\infty
\]
である．したがって任意の $M>0$ に対して，ある $n$ が存在して
\[
\lambda(A\cap Q_n)>M
\]
となる．$A\cap Q_n$ は有界なLebesgue可測集合なので，
有限測度の場合を $A\cap Q_n$ に適用でき，コンパクト集合
$K\subset A\cap Q_n\subset A$ で
\[
\lambda(K)>M
\]
となるものが取れる．よって
\[
\sup\{\lambda(K)\mid K\subset A,\ K\text{ はコンパクト}\}=\infty=\lambda(A)
\]
である．
\end{proof}

\section{まとめ}

本章では，有界な窓で切ったときの内測度と外測度の一致を出発点として，
Lebesgue可測集合とLebesgue測度の基本性質を確立した．
\begin{itemize}
    \item Lebesgue可測集合全体 $\calL_d$ は $\RR^d$ 上の完全加法族，
    すなわち $\sigma$-加法族である
    （\cref{thm:lebesgue-measurable-sigma-algebra}）．
    \item Lebesgue測度 $\lambda:\calL_d\to[0,\infty]$ は
    $\lambda(\emptyset)=0$ と完全加法性を満たす．
    したがって $\lambda$ は $(\RR^d,\calL_d)$ 上の完全加法的測度である
    （\cref{prop:lebesgue-measure-nonnegative,thm:lebesgue-measure-countable-additivity}）．
    \item Lebesgue測度は単調列に関して連続である．すなわち，
    下から増大する可測集合列では和集合の測度が測度の極限に等しく，
    上から減少する可測集合列では最初の集合の測度が有限であれば
    共通部分の測度が測度の極限に等しい
    （\cref{thm:lebesgue-measure-continuity-from-below,thm:lebesgue-measure-continuity-from-above}）．
    \item Lebesgue測度は平行移動不変である．すなわち
    $A\in\calL_d$ と $c\in\RR^d$ に対して $A+c\in\calL_d$ であり，
    $\lambda(A+c)=\lambda(A)$ である
    （\cref{prop:lebesgue-measure-translation-invariant}）．
    \item Lebesgue測度は完備である．実際，$\lambda(N)=0$ となる
    $N\in\calL_d$ と $E\subset N$ に対して，
    $\lambda^*(E)\le\lambda^*(N)=0$ だから $E$ はLebesgue零集合であり，
    $E\in\calL_d$ かつ $\lambda(E)=0$ である
    （\cref{thm:lebesgue-null-measurable}）．
    \item Lebesgue測度は正則である．すなわち任意の $A\in\calL_d$ に対して，
    開集合で外側から，コンパクト集合で内側から測度を近似できる
    （\cref{thm:lebesgue-measure-outer-regularity,thm:lebesgue-measure-inner-regularity}）．
\end{itemize}
